\subsection{Vấn đề của khách hàng mà sản phẩm giải quyết}

Các trung tâm và đơn vị bán khóa học thường có thông tin tương đối chuẩn hóa,
nhưng vẫn phải lặp lại cùng một quy trình sản xuất nội dung và tư vấn ở nhiều
phiên livestream. 1Stream tập trung giải quyết các vấn đề sau:

\begin{itemize}
    \item \textbf{Thiếu người đứng live và giới hạn thời lượng hiện diện:} đội
    ngũ 3--5 người phải đồng thời làm marketing, chuẩn bị nội dung, trực phiên và
    tư vấn tuyển sinh, nên khó duy trì thêm các khung giờ ngoài lịch live thật.
    \item \textbf{Nội dung lặp lại nhưng tốn công sản xuất:} học phí, lịch khai
    giảng, lộ trình học, trình độ đầu vào và ưu đãi thay đổi theo từng chương
    trình; mỗi lần lên sóng vẫn cần tổng hợp lại kịch bản, giọng đọc, hình ảnh và
    phụ đề.
    \item \textbf{Đội tư vấn phải trả lời FAQ nhiều lần:} phần lớn tương tác ban
    đầu xoay quanh một tập câu hỏi quen thuộc. Trả lời thủ công dễ chậm, thiếu
    nhất quán và làm nhân viên ít thời gian hơn cho các trường hợp cần tư vấn
    chuyên sâu.
    \item \textbf{Quy trình livestream bị chia nhỏ:} dữ liệu khóa học, tài sản
    truyền thông, lịch phát, trạng thái phiên và danh sách lead thường nằm ở
    nhiều công cụ khác nhau, gây khó kiểm soát và khó đánh giá hiệu quả theo
    từng đợt.
    \item \textbf{Khó thử nghiệm với ngân sách nhỏ:} thuê host, quay dựng hoặc
    mua một hệ thống dài hạn đều tạo cam kết chi phí trước khi trung tâm biết mô
    hình livestream AI có phù hợp hay không.
    \item \textbf{Rủi ro thông tin và chính sách nền tảng:} câu trả lời sai về
    học phí, lịch học hoặc cam kết đầu ra có thể ảnh hưởng uy tín. Đồng thời,
    tài khoản và nội dung phải đáp ứng điều kiện livestream, bản quyền và yêu
    cầu ghi nhãn nội dung ghi hình trước của từng nền tảng.
\end{itemize}
